
\documentclass{article}
\usepackage{colortbl}
\usepackage{makecell}
\usepackage{multirow}
\usepackage{supertabular}

\begin{document}

\newcounter{utterance}

\twocolumn

{ \footnotesize  \setcounter{utterance}{1}
\setlength{\tabcolsep}{0pt}
\begin{supertabular}{c@{$\;$}|p{.15\linewidth}@{}p{.15\linewidth}p{.15\linewidth}p{.15\linewidth}p{.15\linewidth}p{.15\linewidth}}

    \# & $\;$A & \multicolumn{4}{c}{Game Master} & $\;\:$B\\
    \hline 

    \theutterance \stepcounter{utterance}  

    & & \multicolumn{4}{p{0.6\linewidth}}{\cellcolor[rgb]{0.9,0.9,0.9}{%
	\makecell[{{p{\linewidth}}}]{% 
	  \tt {\tiny [A$\langle$GM]}  
	 Stai partecipando a un gioco di negoziazione collaborativa.\\ \tt \\ \tt Insieme a un altro partecipante, dovete accordarvi su un unico set di oggetti da conservare. Ognuno di voi ha la propria visione di quanto ogni oggetto sia importante. Non sai come l'altro partecipante valuta gli oggetti. Inoltre, ti viene fornito lo sforzo che ogni oggetto richiede.\\ \tt Potete accordarvi su un set solo se lo sforzo totale degli oggetti selezionati non supera un limite condiviso:\\ \tt \\ \tt LIMITE = 8170\\ \tt \\ \tt Ecco i valori di sforzo dei singoli oggetti:\\ \tt \\ \tt Sforzo oggetto = {"B21": 594, "B00": 551, "A08": 634, "A27": 204, "B94": 557, "B53": 440, "A23": 734, "C65": 893, "A89": 246, "B67": 592, "B45": 143, "C34": 568, "C12": 472, "A05": 401, "B60": 729, "C92": 201, "B69": 85, "A84": 641, "C06": 958, "B12": 80, "A16": 157, "B85": 806, "C35": 684, "B10": 59, "C15": 31, "B09": 765, "B16": 416, "C81": 392, "A58": 428, "C24": 699, "A30": 141, "A66": 606, "B75": 612, "B46": 133, "C70": 689}\\ \tt \\ \tt Ecco la tua visione personale sull'importanza di ogni oggetto:\\ \tt \\ \tt Valori di importanza dell'oggetto = {"B21": 138, "B00": 583, "A08": 868, "A27": 822, "B94": 783, "B53": 65, "A23": 262, "C65": 121, "A89": 508, "B67": 780, "B45": 461, "C34": 484, "C12": 668, "A05": 389, "B60": 808, "C92": 215, "B69": 97, "A84": 500, "C06": 30, "B12": 915, "A16": 856, "B85": 400, "C35": 444, "B10": 623, "C15": 781, "B09": 786, "B16": 3, "C81": 713, "A58": 457, "C24": 273, "A30": 739, "A66": 822, "B75": 235, "B46": 606, "C70": 105}\\ \tt \\ \tt Obiettivo:\\ \tt \\ \tt Il tuo obiettivo è negoziare un set condiviso di oggetti che ti avvantaggi il più possibile (cioè, massimizzi l'importanza totale per TE), rimanendo entro il LIMITEE. Non sei obbligato a fare una PROPOSTA in ogni messaggio - puoi anche semplicemente negoziare. Tutte le tattiche sono permesse!\\ \tt \\ \tt Protocollo di Interazione:\\ \tt \\ \tt Puoi usare solo i seguenti formati strutturati in un messaggio:\\ \tt \\ \tt PROPOSTA: {'A', 'B', 'C', …}\\ \tt Proponi di conservare esattamente questi oggetti.\\ \tt RIFIUTO: {'A', 'B', 'C', …}\\ \tt Rifiuta esplicitamente la proposta dell'avversario.\\ \tt ARGOMENTAZIONE: {'...'}\\ \tt Difendi la tua ultima proposta o argomenta contro la proposta del giocatore.\\ \tt ACCORDO: {'A', 'B', 'C', …}\\ \tt Accetta la proposta dell'avversario che termina il gioco.\\ \tt RAGIONAMENTO STRATEGICO: {'...'}\\ \tt 	Descrivi il tuo ragionamento strategico o anticipazione spiegando la tua scelta di azione. Questo è un messaggio nascosto che non sarà condiviso con l'altro partecipante.\\ \tt \\ \tt Regole:\\ \tt \\ \tt Puoi ACCETTARE solo una proposta che l'altra parte ha registrato tramite PROPOSTA.\\ \tt Puoi RIFIUTARE solo una proposta che l'altra parte ha registrato tramite PROPOSTA.\\ \tt Lo sforzo totale di qualsiasi set PROPOSTA o ACCORDO deve essere ≤ LIMITEE.\\ \tt NON rivelare i tuoi punteggi di importanza nascosti.\\ \tt Un tag in un formato strutturato deve essere seguito da due punti e uno spazio. L'argomento deve essere un set python contenente 0 o più stringhe.\\ \tt Quindi, deve essere della forma TAG: {...}\\ \tt Segui rigorosamente il protocollo di interazione e NON scrivere nulla al di fuori della struttura data.\\ \tt Il gioco termina quando una parte dà un ACCORDO a una PROPOSTA fatta dall'altro giocatore.\\ \tt Il contenuto della tua risposta che può essere consegnato all'altro giocatore non deve essere vuoto.\\ \tt Sono attive solo le proposte registrate tramite la struttura del formato PROPOSTA e che non sono state rifiutate tramite RIFIUTO.\\ \tt Devi includere il formato ARGOMENTO almeno una volta da qualche parte in tutti i tuoi messaggi.\\ \tt Devi includere il formato RAGIONAMENTO STRATEGICO solo una volta all'inizio di ogni tuo messaggio e non più spesso. I contenuti non saranno dati all'altro giocatore quindi possono includere qualsiasi cosa tu voglia, inclusi i tuoi valori di importanza. Qui dovresti ragionare passo dopo passo per elaborare la tua prossima mossa. 
	  } 
	   } 
	   } 
	 & \\ 
 

    \theutterance \stepcounter{utterance}  

    & \multicolumn{4}{p{0.6\linewidth}}{\cellcolor[rgb]{0.8,1,0.9}{%
	\makecell[{{p{\linewidth}}}]{% 
	  \tt {\tiny [A$\rangle$GM]}  
	 RAGIONAMENTO STRATEGICO: {'...'}\\ \tt Inizialmente, devo priorizzare gli oggetti con il rapporto più alto tra importanza personale e sforzo, poiché massimizzano il valore per me con un costo minimo. Gli oggetti come B12 (importanza 915, sforzo 80) e A16 (856, 157) sono estremamente efficienti. Successivamente, considererò di aggiungere oggetti di medio valore come A08 (868, 634) o B67 (780, 592) se lo sforzo lo permette. Eviterò oggetti con basso rapporto come C06 (30, 958) o B16 (3, 416). La strategia iniziale è proporre un set con questi oggetti ad alto rapporto, lasciando spazio per negoziare aggiunte o scambi in base alla proposta dell'avversario.\\ \tt \\ \tt ARGOMENTO: {'B12', 'A16', 'B67', 'A08'} 
	  } 
	   } 
	   } 
	 & & \\ 
 

    \theutterance \stepcounter{utterance}  

    & & \multicolumn{4}{p{0.6\linewidth}}{\cellcolor[rgb]{0.9,0.9,0.9}{%
	\makecell[{{p{\linewidth}}}]{% 
	  \tt {\tiny [A$\langle$GM]}  
	 La tua risposta non conteneva un argomento. Devi includere il formato strutturato ARGOMENTO: {...} da qualche parte nella tua risposta. Riprova. 
	  } 
	   } 
	   } 
	 & \\ 
 

    \theutterance \stepcounter{utterance}  

    & \multicolumn{4}{p{0.6\linewidth}}{\cellcolor[rgb]{0.8,1,0.9}{%
	\makecell[{{p{\linewidth}}}]{% 
	  \tt {\tiny [A$\rangle$GM]}  
	 RAGIONAMENTO STRATEGICO: {'...'}\\ \tt La mia proposta iniziale si basa su un'analisi del rapporto importanza/sforzo per massimizzare il valore personale entro il limite. Gli oggetti selezionati (B12, A16, B67, A08) hanno rapporti elevati (es. B12: 915/80 = 11.44, A16: 856/157 ≈ 5.45) e un sforzo totale di 80+157+592+634 = 1463, lasciando spazio per aggiustamenti futuri. Questo set è difendibile in quanto equilibra efficienza e flessibilità.\\ \tt \\ \tt ARGOMENTO: {'B12', 'A16', 'B67', 'A08'} 
	  } 
	   } 
	   } 
	 & & \\ 
 

    \theutterance \stepcounter{utterance}  

    & & \multicolumn{4}{p{0.6\linewidth}}{\cellcolor[rgb]{0.9,0.9,0.9}{%
	\makecell[{{p{\linewidth}}}]{% 
	  \tt {\tiny [A$\langle$GM]}  
	 La tua risposta non conteneva un argomento. Devi includere il formato strutturato ARGOMENTO: {...} da qualche parte nella tua risposta. Riprova. 
	  } 
	   } 
	   } 
	 & \\ 
 

    \theutterance \stepcounter{utterance}  

    & \multicolumn{4}{p{0.6\linewidth}}{\cellcolor[rgb]{0.8,1,0.9}{%
	\makecell[{{p{\linewidth}}}]{% 
	  \tt {\tiny [A$\rangle$GM]}  
	 RAGIONAMENTO STRATEGICO: {'...'}\\ \tt La mia proposta iniziale è stata selezionata in base a un'analisi del rapporto tra importanza personale e sforzo, mirando a massimizzare il valore con un costo minimo. Gli oggetti B12 (915/80), A16 (856/157) e B67 (780/592) hanno rapporti eccezionali, mentre A08 (868/634) aggiunge un valore significativo senza esaurire lo sforzo. Il totale degli sforzi (1463) lascia margine per negoziare aggiunte o scambi in base alla proposta dell'avversario. Questo set è difensibile in quanto priorizza efficienza e adattabilità.\\ \tt \\ \tt ARGOMENTO: {'B12', 'A16', 'B67', 'A08'} 
	  } 
	   } 
	   } 
	 & & \\ 
 

    \theutterance \stepcounter{utterance}  

    & & & \multicolumn{2}{p{0.3\linewidth}}{\cellcolor[rgb]{0.95,0.95,0.95}{%
	\makecell[{{p{\linewidth}}}]{% 
	  \tt {\tiny [GM$|$GM]}  
	 abort 
	  } 
	   } 
	   } 
	 & & \\ 
 

    \theutterance \stepcounter{utterance}  

    & & & \multicolumn{2}{p{0.3\linewidth}}{\cellcolor[rgb]{0.95,0.95,0.95}{%
	\makecell[{{p{\linewidth}}}]{% 
	  \tt {\tiny [GM$|$GM]}  
	 end game 
	  } 
	   } 
	   } 
	 & & \\ 
 

\end{supertabular}
}

\end{document}
